\documentclass[14pt]{extarticle}
\usepackage{fontspec}
\setmainfont{Times New Roman}
\usepackage[a4paper, top=1cm, left=1cm, right=1cm, bottom=2cm]{geometry}
\usepackage{lib}
\usepackage{tikz}
\madethi{111}
\setbool{isTeacher}{false}
\begin{document}
\thongtintruong{CUỐI}{II}{2025-2026}{11}
\begin{khuvuc}{\textbf{Câu trắc nghiệm nhiều phương án lựa chọn.} Mỗi câu hỏi thí sinh chỉ chọn một phương án.}
    \CauTracNghiem
    {Rút gọn biểu thức $P=x^\frac{1}{2}.x^\frac{1}{3}.\sqrt[6]{x}$}
    {$P=\sqrt{x}$}
    {\Dung $P=x$}
    {$P=x\frac{11}{6}$}
    {$P=x\frac{7}{6}$}
    \CauTracNghiem
    {Tập xác định của hàm số $log_{2026}{(x-1)}$ là:}
    {$\left(2;+\infty \right)$}
    {$\left(-\infty;+\infty \right)$}
    {\Dung $\left(1;+\infty \right)$}
    {$\left(-\infty;1 \right)$}
    \CauTracNghiem
    {Nghiệm của phương trình $3^{x-1}=9$ là:}
    {$x=2$}
    {\Dung $x=3$}
    {$x=1$}
    {$x=4$}
    \CauTracNghiem
    {Với $a, b$ là các số thực dương tùy ý và $a \neq 1$, giá trị của $\log_a(a^3)$ bằng:}
    {\Dung $3$}
    {$\dfrac{1}{3}$}
    {$a^3$}
    {$1$}
    \CauTracNghiem
    {Nghiệm của bất phương trình $2^x > 8$ là:}
    {$x>3$}
    {\Dung $x>3$}
    {$x>4$}
    {$x<3$}
    \CauTracNghiem
    {Đạo hàm của hàm số $y=x^4-3x^2+1$ là:}
    {$y'=4x^3-3x$}
    {$y'=x^3-6x$}
    {\Dung $y'=4x^3-6x$}
    {$y'=4x^3-6x+1$}
    \CauTracNghiem
    {Cho hàm số $y=sin(x)$. Đạo hàm y' bằng: }
    {\Dung $y'=cos(x)$}
    {$y'=x.cos(x)$}
    {$y=-cos(x)$}
    {$y'=-sin(x)$}
    \CauTracNghiem
    {Hàm số $y=\dfrac{x+1}{x-2}$ có đạo hàm là: }
    {\Dung $y'=\dfrac{3}{(x-2)^2}$}
    {$y'=\dfrac{-3}{(x-2)^2}$}
    {$y'=\dfrac{-1}{(x-2)^2}$}
    {$y'=\dfrac{1}{(x-2)^2}$}
    \CauTracNghiem
    {Trong không gian, cho hai đường thẳng $a$ và $b$ song song. Nếu đường thẳng $c$ vuông góc với $a$ thì:}
    {$c // b$}
    {\Dung $c \perp b$}
    {$c \equiv b$}
    {$c$ và $b$ chéo nhau}
    \CauTracNghiem
    {Cho đường thẳng $d$ vuông góc với mặt phẳng $(P)$. Đường thẳng $d$ sẽ vuông góc với:}
    {\Dung Mọi đường thẳng nằm trong $(P)$}
    {Chỉ một đường thẳng nằm trong $(P)$}
    {Chỉ các đường thẳng đi qua giao điểm của $d$ và $(P)$}
    {Không có đường thẳng nào trong $(P)$}
    \CauTracNghiem
    {Cho hình chóp $S.ABC$ có $SA \perp (ABC)$. Khi đó góc giữa $SB$ và mặt phẳng $(ABC)$ là góc:}
    {$\widehat{SBC}$}
    {$\widehat{ASB}$}
    {\Dung $\widehat{SBA}$}
    {$\widehat{SCA}$}
    \CauTracNghiem
    {Cho hai mặt phẳng $(P)$ và $(Q)$ vuông góc với nhau theo giao tuyến $d$. 
    Nếu đường thẳng $a$ nằm trong $(P)$ và $a \perp d$ thì:}
    {\Dung $a \perp (Q)$}
    {$a // (Q)$}
    {$a \subset (Q)$}
    {$a$ cắt $(Q)$ nhưng không vuông góc}
    \end{khuvuc}
    \begin{khuvuc}{\textbf{Câu trắc nghiệm đúng sai.} Thí sinh trả lời từ câu 1 đến câu 4. Trong mỗi ý a), b), c), d) ở mỗi câu, thí sinh chọn đúng hoặc sai.}
        \CauDungSai
        {Cho biểu thức $P(x) = \left( x^{\frac{2}{3}} \right)^3$ và hàm số $f(x) = 2^x - 4$.}
        {\Dung Tập nghiệm của bất phương trình $f(x) > 0$ là $(2; +\infty)$.}
        {\Dung Phương trình $f(x) = 0$ có nghiệm duy nhất $x = 2$. }
        {\Dung Với $x > 0$, biểu thức $P(x)$ được rút gọn thành $x^2$.}
        {\Sai Hàm số $y = f(x)$ là hàm số nghịch biến trên $\mathbb{R}$.}
        \CauDungSai
        {Cho hàm số $g(x) = \log_2(x^2 - 3x)$.}
        {\Sai Tập xác định của hàm số $g(x)$ là $D = (0; 3)$. }
        {\Dung Bất phương trình $g(x) \le 2$ có nghiệm nguyên dương lớn nhất là $x = 4$.}
        {\Dung Giá trị của biểu thức $g(4) = 2$.}
        {\Dung Phương trình $g(x) = 1$ có tổng các nghiệm bằng 3.}
        \\
        \CauDungSai
        {Cho hàm số $y = \dfrac{x + 2}{x - 1}$ có đồ thị là $(C)$.}
        {\Sai Đạo hàm của hàm số đã cho là $y' = \frac{-1}{(x - 1)^2}$.}
        {\Dung Hệ số góc của tiếp tuyến của $(C)$ tại điểm có hoành độ $x = 2$ là $k = -3$.}
        {\Dung Phương trình tiếp tuyến của $(C)$ tại điểm $M(2; 4)$ là $y = -3x + 10$. }
        {\Sai Đạo hàm của hàm số tại $x = 0$ có giá trị dương.}
        \CauDungSai
        {Cho hình chóp $S.ABCD$ có đáy $ABCD$ là hình vuông tâm $O$ cạnh $a$. 
        Cạnh bên $SA = a\sqrt{3}$ và $SA \perp (ABCD)$.\\
        \centerline{
            \begin{tikzpicture}[scale=0.5, font=\footnotesize, line join=round, line cap=round]
                \coordinate (A) at (0,0);
                \coordinate (B) at (-2,-3);
                \coordinate (D) at (8,0);
                \coordinate (C) at ($(B)+(D)-(A)$);
                \coordinate (S) at (0,5);
                \coordinate (O) at (intersection of A--C and B--D);

                % Vẽ các cạnh khuất (dashed)
                \draw[dashed] (S)--(A) (A)--(B) (A)--(D) (A)--(C) (B)--(D);
                
                % Vẽ các cạnh thấy (liền nét)
                \draw[thick] (S)--(B)--(C)--(D)--cycle;
                \draw[thick] (S)--(C);

                % Ký hiệu vuông góc tại A
                \draw (0.2,0)--++(0,0.2)--++(-0.2,0); % SA vuông góc AD

                % Dán nhãn các điểm
                \node[left] at (A) {$A$};
                \node[below left] at (B) {$B$};
                \node[below right] at (C) {$C$};
                \node[right] at (D) {$D$};
                \node[above] at (S) {$S$};
                \node[below] at (O) {$O$};

                % Vẽ điểm O
                \fill (O) circle (0.5pt);
        \end{tikzpicture}
        }
        \vspace{-10pt}
        }
        {\Dung Đường thẳng $BD$ vuông góc với mặt phẳng $(SAC)$. }
        {\Sai Góc giữa đường thẳng $SC$ và mặt phẳng đáy $(ABCD)$ là góc $\widehat{SAC}$. }
        {\Dung Khoảng cách từ điểm $S$ đến mặt phẳng $(ABCD)$ bằng $a\sqrt{3}$.}
        {\Sai Thể tích khối chóp $S.ABCD$ là $V = a^3\sqrt{3}$.}
    \end{khuvuc}
    \begin{khuvuc}{\textbf{Câu trắc nghiệm trả lời ngắn.} Thí sinh trả lời từ câu 1 đến câu 6.}
        \CauTuLuan{Cho biểu thức $P = \frac{x \cdot \sqrt[3]{x^2}}{\sqrt{x}}$ với $x > 0$. 
        Khi rút gọn, ta được $P = x^{\frac{a}{b}}$ (với $\frac{a}{b}$ là phân số tối giản). 
        Tính giá trị của $T = a + b$.}
        {
            
        }{}
        \CauTuLuan{Tìm tập nghiệm của bất phương trình $\left( \dfrac{1}{2} \right)^{x^2 - 3x} \ge 4$. 
        Tập nghiệm có dạng $[m; n]$. Tính giá trị $S = m + n$.}
        {
            
        }{}
        \CauTuLuan{Cho $\log_2 5 = a$ và $\log_2 3 = b$. Tính giá trị của biểu thức $M = \log_3 135$ theo $a$ và $b$. 
        Nếu $M = \dfrac{ma + nb}{kb}$, hãy tính tổng $m + n + k$.}
        {
            
        }{}
        \CauTuLuan{Cho hàm số $f(x) = \dfrac{1}{3}x^3 - 2x^2 + 3x + 2026$. 
        Tính giá trị của đạo hàm tại điểm $x = 1$.}
        {
            
        }{}
        \CauTuLuan{Một đoàn tàu chuyển động với phương trình $s(t) = 10t - 0,5t^2$ ($t$ tính bằng giây, $s$ tính bằng mét). 
        Tính vận tốc của đoàn tàu (m/s) tại thời điểm nó đi được $5$ giây kể từ lúc bắt đầu chuyển động.}
        {
            
        }{}
        \CauTuLuan{Một kiến trúc sư thiết kế mái nhà hình chóp cụt cho một biệt thự. 
        Phần mái cao phía trên là hình chóp đều $S.ABCD$ có đáy $ABCD$ là hình vuông cạnh $6\text{m}$. 
        Sau khi thi công, do yêu cầu về thoát nước mưa, kiến trúc sư cần tính toán góc giữa mặt bên (mái dốc) và mặt đáy (trần nhà). Biết rằng chiều cao của khối chóp mái là $SO = 3\sqrt{3}\text{m}$ (với $O$ là tâm của đáy $ABCD$). 
        Tính số đo của góc tạo bởi mặt bên và mặt phẳng đáy (đơn vị: độ).}
        {
            
        }{}
    \end{khuvuc}
\cuoidethi
\end{document}